\subsection{The expansion of the master equation}
\label{App:LargeN-expansion}
Inherent stochasticity ---as opposed to externally-induced fluctuations--- is fully described by the master equation. In one-dimensional, one-step processes, this equation describes a time-dependent, integer-valued stochastic variable, $n$, through the temporal evolution of its associated probability distribution, $P(n,t)$. We know that, in general, deterministic descriptions are more accurate the larger a system is. This means that the relative importance of the fluctuations around the deterministic average decreases as the system becomes larger and larger. This systematic expansion furnishes a general approximation to the master equation that highlight this intuitive idea. The method gives rise to an equation for the average of the variable, $n$, ---a.k.a {\em macroscopic law} \cite{vanKampen92}---, and an equation for the fluctuations around this average, which becomes more and more accurate as more and more terms, weighted by decreasing, negative powers of a {\em size} parameter, are considered. This {\em size} parameter will depend on the system under study. 
\smallskip

We will present this method through a biological example. Most marine invertebrate species have a planktonic larva. This stage can be seen as a dispersal strategy. However, adult life is bentonic. After some time in the plankton, larvae settle down and a juvenile life stage starts. During this juvenile stage individuals mature, find mates, and reproduce, sending again the offspring back to the plankton. We will consider the dynamics of a juvenile population of a given marine species in an area of size $\Omega$. Juveniles enter the population as a constant influx from the plankton. They keep this juvenile state until finding a mate. So, pairwise encounters determine the rate at which juveniles disappear from the area. Here we are not considering additional juvenile mortality. This dynamics can be summarized by the following transition rates. The probability of a new juvenile entering the area will be proportional to that area $\Omega$. Apart from an intrinsic probability controlling the mating process, the probability of disappearance of a juvenile will be proportional to the probability of a pairwise encounter, and this probability will be proportional to the inverse of the area, $\Omega$. 
Therefore, we have:
\begin{eqnarray}
	\label{g_n_exp}
	g_n &=& k_{1}\,\Omega\\
	r_n &=& k_{2}\,\frac{n(n-1)}{\Omega},
	\label{r_n_exp}
\end{eqnarray}
where $k_{1}$ is the influx rate per unit area and $k_{2}$ controls the intrinsic mating rate once a pair-wise encounter has taken place. As usual, these rates enable the writing of an equation for the temporal evolution of the number of juveniles in the area $\Omega$, $P(n,t)$, which is, as we know, the master equation, Eq. (\ref{me_1}). In this context, it will be very useful to introduce two operators that simplify this equation:
\begin{eqnarray}
	\label{E_1}
	\textup{E}f(n)      &=& f(n+1)\\
	\textup{E}^{-1}f(n) &=& f(n-1)
	\label{E_2}
\end{eqnarray}
Notice that the above definitions transform the Eq. (\ref{me_1}) into:
\begin{equation}
	\label{me_E}
	\frac{dP(n,t)}{dt}=(\textup{E}-1)r_n\,P(n,t) + (\textup{E}^{-1}-1)g_n\,P(n,t)
\end{equation}  
By using Eqs. (\ref{g_n_exp})-(\ref{r_n_exp}), and also a shorter notation for 
$P(n,t)\equiv\,p_n$, we have:
\begin{equation}
	\label{me_E}
	\frac{dp_n}{dt}=\frac{k_2}{\Omega}(\textup{E}-1)n(n-1)\,p_n + k_{1}\Omega(\textup{E}^{-1}-1)\,p_n\,.
\end{equation}  
which is valid for $n>0$. For $n=0$ only the arrival of new juveniles is possible:
\begin{equation}
	\frac{dp_0}{dt}= -k_{1}\Omega\,p_0.
\end{equation}
\smallskip

The expansion of the master equation is based on the following assumption. The number of juveniles at some time $t$ can be approximated by:
\begin{equation}
	\label{central}
	n = \Phi(t)\,\Omega + \Omega^{\frac{1}{2}}\,\xi
\end{equation}
This assumes that the temporal evolution of the variable $n$ can be described by a combination of the a new variable $\Phi(t)$ which entails the temporal evolution of the mean value of $n$, and a new variable $\xi$, which controls for the importance of fluctuations around this mean value. Notice that the relative importance of fluctuations around mean densities scale with the size of the area $\Omega^{-1/2}$, according to the general intuition that fluctuations are relatively more important in small systems, i.e., at low numbers. More precisely, notice that:
\begin{enumerate}
	\item The density of individuals follows the evolution of the mean density $\Phi(t)$ as the size of the system is large:
	\begin{equation}
		\frac{n}{\Omega}=\Phi(t)+\xi\,\Omega^{-\frac{1}{2}}\, \Rightarrow\, \lim_{\Omega \to \infty}\frac{n}{\Omega}=\Phi(t)\,.
	\end{equation}    
	\item The expected value for the number of juveniles should be written as $\Phi(t)\,\Omega$, which requires the expected value for the stochastic variable $\xi$ to be zero:
	\begin{equation}
		<n> = \Phi(t)\,\Omega \Rightarrow <\xi>=0\,.
	\end{equation} 
	\item The variance in the number of juveniles increases with $\Omega$, but the relative importance of fluctuations in relation to the size of the system will vanishes \vicente{typo} as $\Omega$ tends to infinity:
	\begin{eqnarray}
		\textup{Var}(n)&=&<(n-<n>)^2>=<(\xi\,\Omega^{\frac{1}{2}})^2> = \Omega\,<\xi^2>=\Omega\,\textup{Var}(\xi)\,\nonumber\\
		&\Rightarrow& \textup{Var}(\xi)=\frac{\textup{Var}(n)}{\Omega}\,\to\,0\,\,\textup{as}\,\,\Omega\,\to\,\infty\,. 
	\end{eqnarray}
\end{enumerate}

Our interest now is to translate the master equation for $P(n,t)$ (Eq. \ref{me_E}) into an equivalent equation for a probability distribution of the new stochastic variable, $\xi$, through the suggested change of variables ---see Eq. (\ref{central}). The goal is obtaining, first, as the leading contribution of the expansion, an equation for the temporal evolution of mean density values, and, second, an equation for the temporal evolution of fluctuations, which include further terms in decreasing, negative powers of the {\em size} parameter, $\Omega$. In oder to make this transformation, first, we note that:
\begin{enumerate}
	\item The probability distribution $P(n,t)$ can be written in terms of $\xi$ by introducing the central expression for $n$ (Eq. \ref{central}):
	\begin{equation}
		\label{eq:Pi_Def}
		p_n \equiv P(n,t) = P(\Phi(t)\,\Omega + \Omega^{\frac{1}{2}}\,\xi,t)\,\equiv\,\Pi(\xi,t)
	\end{equation}
	Here, we will take Eq. (\ref{eq:Pi_Def}) as the definition of the function $\Pi(\xi,t)$\footnote{However, $P(n,t)$ and $\Pi(\xi,t)$ are two probability distributions whose associated stochastic variables are related to each other by a linear change ($n~=~\Phi\Omega+\Omega^{1/2}\xi$). This will introduce a multiplicative factor in the above given definition: $\Pi(\xi,t)~=~\Omega^{1/2}\,P(\Phi\Omega+\Omega^{1/2}\xi,\,t)$, but this will not alter alter the rest of this derivation in any meaningful way. Notice that the set of ODEs in Eq. (\ref{me_E}) determine the probabilities except a multiplicative constant.}.     
	\item The time derivatives arising in the master equation (\ref{me_E}) are performed by keeping $n$ constant for each equation of the set. We must bear this in mind, when we express these time derivatives in terms of the new variable $\xi$ and the new probability distribution $\Pi(\xi,t)$. More precisely, since $n=\Phi(t)\,\Omega + \Omega^{\frac{1}{2}}\,\xi=\textup{Constant}$, if we derive this expression with respect to time, we will have :
	\begin{eqnarray}
		\label{time_derivative}
		\Omega\,\frac{d\Phi}{dt}+\Omega^{\frac{1}{2}}\,\frac{d\xi}{dt} &=& 0\nonumber\\
		&\Rightarrow&\,\frac{d\xi}{dt}=-\Omega^{\frac{1}{2}}\,\frac{d\Phi}{dt}\,.
	\end{eqnarray}
	\item The new probability distribution $\Pi(\xi,t)$ has an explicit and an implicit time dependence through the variable $\xi$. The last consideration allows to calculate these total time derivative as:
	\begin{equation}
		\label{time_derivative_2}
		\frac{d\Pi(\xi,t)}{dt}=\frac{\partial\Pi}{\partial\,t} + \frac{\partial\Pi}{\partial\,\xi}\frac{d\xi}{dt} 
	\end{equation}
	which can be written, by using Eq. (\ref{time_derivative}), as:
	\begin{equation}
		\frac{d\Pi(\xi,t)}{dt}=\frac{\partial\Pi}{\partial\,t} - \Omega^{\frac{1}{2}}\,\frac{d\Phi}{dt} \frac{\partial\Pi}{\partial\,\xi}	
		\label{eq:dPi}
	\end{equation}
\end{enumerate}
\smallskip

Since $\Pi(\xi,t)$ has been defined as $P(\Phi(t)\,\Omega + \Omega^{\frac{1}{2}}\,\xi,t)$, Eq. (\ref{eq:dPi}) completes the transformation of the left hand side of the master equation ---see Eq. (\ref{me_E}). However, we still need to convert the operators $\textup{E}$ and $\textup{E}^{-1}$ and the transition rates into functions of $\xi$. By using Eq. (\ref{central}), transition rates given by Eqs. (\ref{g_n_exp})-(\ref{r_n_exp}) can be readily written in terms of $\xi$. After some algebra, we obtain:\vicente{The first term of Eq 50 should have Phi squared}
\begin{eqnarray}
	\label{g_n_Exp}
	g_n & = & k_{1}\,\Omega\\
	r_n & = & k_{2}\,\Omega\,\Phi+2k_{2}\,\xi\,\Phi\Omega^{\frac{1}{2}}-k_{2}(\Phi-\xi^2) - k_{2}\,\xi\,\Omega^{-\frac{1}{2}}\,. 
	\label{r_n_Exp}
\end{eqnarray}
Now, we should investigate how the operators $\textup{E}$ and $\textup{E}^{-1}$, defined by Eqs. (\ref{E_1})-(\ref{E_2}), change a function that depends on $\xi$ rather than on $n$. This is quite easy if we note that, by using Eq. (\ref{central}):
\begin{equation}
	n+1 = \Phi\,\Omega + \Omega^{\frac{1}{2}}\,(\xi+\Omega^{-\frac{1}{2}})
\end{equation}
Therefore, the action of the operators $\textup{E}$ and $\textup{E}^{-1}$ on functions depending on $\xi$ will be:
\begin{eqnarray}
	\textup{E}\,f(n) &=& \textup{E}\,F(\xi) = F(\xi + \Omega^{-\frac{1}{2}})\\
	\textup{E}^{-1}\,f(n) &=& \textup{E}^{-1}\,F(\xi) = F(\xi - \Omega^{-\frac{1}{2}})
\end{eqnarray} 
The key point of the systematic expansion of the master equation is to write these operators as a Taylor expansion around $\xi$, where the accuracy of the expansion improves as the size of the system, $\Omega$, tends to infinity. Alternatively, as the size of the systems shrinks, in order to keep system description accurate, we will need more and more terms in the expansion. Notice that the Taylor expansion is left in terms of negative powers of $\Omega$:
\begin{eqnarray}
	\label{E_3}
	\textup{E}\,f(n) & = &\textup{E}\,F(\xi) = F(\xi)+\Omega^{-\frac{1}{2}}\frac{dF}{d\xi}+\frac{1}{2}\Omega^{-1}\frac{d^2F}{d\xi^2}+\ldots\\
	\textup{E}^{-1}\,f(n) & = &\textup{E}^{-1}\,F(\xi) = F(\xi)-\Omega^{-\frac{1}{2}}\frac{dF}{d\xi}+\frac{1}{2}\Omega^{-1}\frac{d^2F}{d\xi^2}-\ldots
	\label{E_4}
\end{eqnarray} 
Therefore, in general, we can write the operators arising in the master equation, Eq. (\ref{me_E}), as true operators in terms of negative powers of $\Omega$:
\begin{eqnarray}
	\label{E_5}
	\textup{E}\,-1      & = &  \Omega^{-\frac{1}{2}}\frac{\partial}{\partial\xi}+\frac{1}{2}\Omega^{-1}\frac{\partial^2}{\partial\xi^2}+\ldots\\
	\textup{E}^{-1}\,-1 & = & -\Omega^{-\frac{1}{2}}\frac{\partial}{\partial\xi}+\frac{1}{2}\Omega^{-1}\frac{\partial^2}{\partial\xi^2}-\ldots\,,
	\label{E_6}
\end{eqnarray}
where we have used partial derivatives to emphasize the fact that these operators act on functions in general dependent on both $\xi$ and time. 

Finally, we have all the necessary ingredients to fully transform Eq. (\ref{me_E}) in terms of the new variable $\xi$ and its associated probability distribution $\Pi(\xi,t)$. First, to keep generality, let us only make use of Eqs. (\ref{eq:dPi}), and the expressions for the operators (\ref{E_5})-(\ref{E_6}), and introduce them into the master equation, Eq. (\ref{me_E}). This is what we obtain:\vicente{And here we are missing the 1/2 in the second derivatives}
\begin{eqnarray}
	\frac{\partial\Pi}{\partial\,t}-\Omega^{\frac{1}{2}}\,\frac{d\Phi}{dt}\frac{\partial\Pi}{\partial\,\xi} & = & \left\{\Omega^{-\frac{1}{2}}\partial_{\xi}+\Omega^{-1}\partial_{\xi}^2+\dots\right\}(r_n\,\Pi(\xi,\,t))\nonumber\\
	&   & \left\{-\Omega^{-\frac{1}{2}}\partial_{\xi}+\Omega^{-1}\partial_{\xi}^2-\dots\right\}(g_n\,\Pi(\xi,\,t))\,,
\end{eqnarray}
which can be written as:
\begin{eqnarray}
	\label{Expan}
	\frac{\partial\Pi}{\partial\,t}-\Omega^{\frac{1}{2}}\,\frac{d\Phi}{dt}\frac{\partial\Pi}{\partial\,\xi} & = &\Omega^{-\frac{1}{2}}\partial_{\xi}\left[(r_n-g_n)\,\Pi\right]\nonumber\\
	&   &\Omega^{-1}\partial_{\xi}^2\left[(r_n+g_n)\,\Pi\right]\nonumber\\
	&   &\dots\,,
\end{eqnarray}
This is still valid for any one-step process, defined by rates $g_n$ and $r_n$. Now, by introducing the transition rates (\ref{g_n_Exp})-(\ref{g_n_Exp}) in Eq. (\ref{Expan}), we can now group terms in powers of $\Omega$. After some algebra, this process give rise to two equations:
\begin{enumerate}
	\item Grouping terms in powers of $\Omega^{1/2}$, we obtain the equation for the temporal evolution of the density of juveniles, $\Phi(t)$, the macroscopic law:
	\begin{equation}
		\label{macroscopic_law}
		\frac{d\Phi}{dt} = k_{1} - k_{2}\,\Phi^2\,.
	\end{equation}
	This deterministic description gives the evolution of the average density of individuals, $\Phi=n/\Omega$ as a result of a constant influx of new juveniles and a density-dependent disappearance consequence of pair-wise encounters. 
	\item Grouping terms in powers $\Omega^0$, a second equation emerges, which governs the fluctuations around the expected value of $n$. It can be grouped, in turn, in terms affected by decreasing, negative powers of $\Omega$. This further terms can be regarded as corrections to the leading order in $\Omega^0$. This equation is:
	\begin{eqnarray}
		\label{linear}
		\frac{\partial\Pi}{\partial\,t}&=&-\frac{d\alpha_1}{d\Phi}\frac{\partial(\xi\Pi)}{\partial\,\xi}+\frac{1}{2}\alpha_2(\Phi)\frac{\partial^2\Pi}{\partial\,\xi^2}\\
		&+&\Omega^{-\frac{1}{2}}\left[\dots\,\,\dots\right]\nonumber\\
		&+&\Omega^{-\frac{3}{2}}\left[\dots\,\,\dots\right]\nonumber
	\end{eqnarray} 
	where
	\begin{eqnarray}
		\label{a_1}
		\alpha_1(\Phi) &\equiv& k_{1} - k_{2}\,\Phi^2 \\ 
		\alpha_2(\Phi) &\equiv& k_{1} + k_{2}\,\Phi^2
		\label{a_2}
	\end{eqnarray} 
\end{enumerate}

If we neglect orders further than $\Omega^{0}$, we obtain the so-called linear noise approximation \cite{vanKampen92}. Within this approximation, the macroscopic law and the equation for the fluctuations can be written in general terms as:
\begin{eqnarray}
	\frac{d\Phi}{dt} &=& \alpha_1(\Phi)\\
	\frac{\partial\Pi}{\partial\,t}&=&-\frac{d\alpha_1}{d\Phi}\frac{\partial(\xi\Pi)}{\partial\,\xi}+\frac{1}{2}\alpha_2(\Phi)\frac{\partial^2\Pi}{\partial\,\xi^2}
	\label{linear_noise}
\end{eqnarray}
Notice that $\alpha_1(\Phi)$ (Eq. \ref{a_1}) and $\alpha_2(\Phi)$ (Eq. \ref{a_2}) only differ in one sign. This is completely general. The function $\alpha_1(\Phi)$ is equal to the right hand side of macroscopic law, while the function $\alpha_2(\Phi)$ contains also the loss and growth terms of the macroscopic law but all with positive signs. 

Eq. (\ref{linear_noise}) is a linear Focker-Planck equation\footnote{The general form of a one-dimensional Fokker-Planck equation is: $$ \frac{\partial\Pi}{\partial\,t} = -\frac{\partial(A(\xi,t)\Pi)}{\partial\,\xi} + \frac{\partial^2 B(x,t)\Pi}{\partial\,\xi^2}$$. If $A(\xi,t)=\xi \, a(t)$ and $ B(\xi,t)=\xi \, b(t)$, we say that the Focker-Planck equation is linear. In our case, $A(\xi,t)=\alpha'_1(\Phi)\,\xi$, and $B(\xi,t) = 1/2\,\alpha_2(\Phi)$, which fulfills linearity. This is why this approximation is called a {\em linear noise approximation.}} whose coefficients depend on time through $\Phi$. The solution can be shown to be a Gaussian:
\begin{equation}
	\label{Pi}
	\Pi(\xi,\,t) \approx \exp\left(-\frac{(\xi-<\xi>)^2}{2\,\textup{Var}(\xi)}\right)\,,
\end{equation}
where $<\xi>$ and $\textup{Var}(\xi)$ are the usual time-dependent Gaussian parameters. The temporal evolution of them can be calculated from Eq. (\ref{linear_noise}) by noticing the following equalities, where an integration by parts has been performed in those cases when was necessary:
\begin{eqnarray}
	\int_{-\infty}^{+\infty}\xi\frac{\partial\Pi}{\partial\,t}d\xi &=& \partial_t<\xi>\\
	\int_{-\infty}^{+\infty}\xi^2\frac{\partial\Pi}{\partial\,t}d\xi &=& \partial_t<\xi^2>\\
	\int_{-\infty}^{+\infty}\xi\frac{\partial(\xi\Pi)}{\partial\,\xi}d\xi &=& -\,<\xi>\\
	\int_{-\infty}^{+\infty}\xi^2\frac{\partial(\xi\Pi)}{\partial\,\xi}d\xi &=& -2\,<\xi^2>\\
	\int_{-\infty}^{+\infty}\xi\frac{\partial^2(\Pi)}{\partial\,\xi^2}d\xi &=& 0\\
	\int_{-\infty}^{+\infty}\xi^2\frac{\partial^2(\Pi)}{\partial\,\xi^2}d\xi &=& 2\,.
\end{eqnarray}
Therefore, we finally end up with two equations:
\begin{eqnarray}
	\label{xi}
	\partial_t<\xi>   & = & \frac{d\alpha_1}{d\Phi}\,<\xi>\\
	\partial_t<\xi^2> & = & 2\,\frac{d\alpha_1}{d\Phi}\,<\xi^2>+\alpha_2    
	\label{xi_2}
\end{eqnarray}
Eq. (\ref{xi}) gives the temporal evolution of expected value of $\xi$. As commented, since fluctuations are required to be centered around the mean value of $n$, consistency also requires that $<\xi>=0$. Furthermore, if our initial condition did not fulfill this requirement, the temporal evolution of $<\xi>$ would tend to make this initially non-centered fluctuation equal to zero, due to Eq. (\ref{xi}). However, if we take $<\xi>_0=0$, then $<\xi>_t=0$ for any further $t$. As as consequence, we will have that $\textup{Var}(\xi)=<\xi^2>$, with Eq. (\ref{xi_2}) governing its temporal evolution. This must be solved together with the macroscopic law to have a complete description of the temporal evolution of the initial probability distribution, $P(n,t)$, within the linear noise approximation. By making use of Eq. (\ref{Pi}) and the change of variable relating $n$ to $\xi$, given by Eq. (\ref{central}), this distribution can be written as:
\begin{equation}
	P(n,t) = \frac{1}{\sqrt{2\pi\,\Omega\,\textup{Var}(\xi)}}\exp\left(-\frac{(n-\Omega\Phi(t))^2}{2\,\Omega\,\textup{Var}(\xi)}\right)
\end{equation} 
where $\Phi(t)$ and $\textup{Var}(\xi)$ are time-dependent functions which give how the Gaussian moves as a whole while its variance changes according to Eq. (\ref{xi_2}).